\section{结论}
\label{sec:conclusion}

本文在受控图分类协议下比较了五种残差路由拓扑，发现峰值准确率、跨数据集一致性和表示连续性之间存在系统性权衡。本文最核心的贡献，是在不改变主干、池化规则和评估协议的前提下，构建出一个能将这些权衡定量呈现出来的比较框架。

回到三个研究问题：（1）最强架构取决于数据集区间，而不是单一普适的残差规则；在面向论文主线的蛋白质导向比较中，NodeResGNN 在 PROTEINS 上最优，而 GraphResGNN 在 DD 上最优。ENZYMES 指向相同的 GraphResGNN 方向，但主要作为支持性压力测试，而非与主证据并列。（2）交叉残差交换有用，但具有选择性：它通常优于 PlainGNN，在 DD 和 Mutagenicity 上保持竞争力，并在 DD 的算子级比较中尤为值得关注，但它并不会稳定超过最优残差基线。（3）路由策略是有意义的次级设计轴：在针对性消融中，轻度稀疏路由最常最优，而可学习稠密路由仍在最强的图级残差目标上保持最佳。

在 PROTEINS（GCNConv）上的深度扩展分析揭示了一个值得进一步验证的“CKA--性能悖论”：表示连续性最高的 NodeResGNN 并未最常获胜，而表示连续性最低的 GraphResGNN 却在更广泛的数据集--算子 sweep 中获胜最多。这一现象提示，但尚未最终证明，层间的结构化表示变化可能是有益而非有害的。

本文最强的结论由以下证据共同支撑：蛋白质导向核心上的分层五折 benchmark、针对性的五折消融以及 $L=1$--$8$ 的深度扩展分析。本文解决的是中等规模生物分子 benchmark 上的受控图分类行为；未来工作应继续检验 CKA--性能悖论，以及“稀疏路由 + 交叉残差”这一组合，是否能推广到更大图、更丰富的生物输入和超越 benchmark 级别的任务场景。
